\chapter{Background}\label{chap:background}
This chapter introduces the core technologies and concepts that support the development of the padel analysis system. The system integrates real-time object detection, motion tracking, and machine learning classification to enable automated sports analytics. The main tools explored include YOLOv8 for object detection, OpenCV with optical flow for ball tracking, and a Random Forest classifier for bounce classification.

\section{Concepts Overview} \label{sec:s1}

\subsection{Artificial Intelligence (AI)}

Artificial Intelligence (AI) studies methods for building systems that perform tasks typically requiring human intelligence, such as learning, reasoning, perception, and decision-making. Advances in computational power, data availability, and algorithms have accelerated AI development and enabled its broad use across domains including healthcare, robotics, autonomous vehicles, and sports analytics. In this work, AI provides the foundational methods—particularly computer vision and machine learning—used to analyze padel videos and generate training feedback.

\subsection{Machine Learning}

Machine Learning (ML) is a subfield of AI that enables systems to learn from data and improve performance without being explicitly programmed for every scenario. ML typically involves a training phase, where a model's parameters are optimized on labeled (or unlabeled) data, and an inference phase, where the trained model is applied to new data. Common paradigms include supervised learning (with labeled data), unsupervised learning (discovering structure in unlabeled data), and reinforcement learning (learning via interaction and rewards). Datasets are partitioned into training, validation, and test sets to evaluate generalization, and optimization algorithms such as gradient descent are used to minimize a loss function.

In the proposed Padel Trainer, supervised deep learning models are trained to detect the ball, players, and court regions from annotated video frames.

\subsection{Deep Learning and Neural Networks}

Deep learning is a class of machine learning that employs artificial neural networks with multiple layers to learn hierarchical representations from raw data. Neural networks consist of interconnected units (neurons) whose outputs are computed as nonlinear functions of weighted sums of their inputs. A typical neuron computes
\[
y = f\Big(\sum_{i=1}^{n} w_i x_i + b\Big),
\]
where \(x_i\) are inputs, \(w_i\) are weights, \(b\) is a bias, and \(f\) is an activation function (e.g., ReLU, sigmoid). Stacking many layers enables models to learn complex, abstract features: early layers capture low-level patterns (edges, textures) while deeper layers capture semantic concepts.

Convolutional Neural Networks (CNNs) are specialized architectures for image data that use convolutional layers to extract spatial features. CNNs form the basis of modern object detection systems used in this thesis, providing automatic feature extraction that outperforms traditional hand-crafted features on large datasets.

\subsection{Computer Vision}

Computer vision focuses on algorithms that enable machines to interpret visual data from images and video. Video analysis extends image methods to exploit temporal information, enabling tasks such as object tracking, motion estimation, and event detection. In sports analytics, computer vision is widely applied for player tracking, ball tracking, and automated event detection; these tasks require real-time or near-real-time processing for practical feedback.

\subsection{Object Detection}

Object detection locates and classifies objects within images by predicting bounding boxes and associated confidence scores. Modern detectors based on CNNs are typically categorized as two-stage or one-stage. Two-stage detectors generate region proposals and then classify them, trading speed for accuracy. One-stage detectors perform localization and classification in a single pass, enabling real-time performance—an important property for sports tracking where objects move rapidly.

This thesis uses object detection to localize the ball, players, and court features in each video frame.

\subsection{YOLO and YOLOv8}

YOLO (You Only Look Once) is a family of one-stage detectors designed for real-time object detection. YOLO divides the image into a grid and predicts bounding boxes, class probabilities, and confidence scores for each grid cell in a single network pass, offering an effective trade-off between speed and accuracy.

YOLOv8 is a modern incarnation offering architectural and training improvements, anchor-free predictions, and variants ranging from Nano (n) to Extra Large (x). For this project, YOLOv8 Nano (YOLOv8n) was chosen for its lightweight architecture and suitability for edge deployment. Custom models were trained for ball, player, and court detection and later exported to ONNX for edge inference.

\subsection{Optical Flow}

Optical flow estimates motion between consecutive frames by computing displacement vectors for pixels or selected feature points. Sparse methods (e.g., Lucas--Kanade) track a set of salient features and are computationally cheaper than dense methods, making them appropriate for real-time tracking. Optical flow is used as a fallback tracker in this thesis when the detector fails temporarily (for example due to motion blur or partial occlusion), helping maintain a continuous trajectory.

\subsection{Kalman Filter}

The Kalman Filter is a recursive estimator for linear dynamical systems that fuses predictions from a motion model with noisy observations to produce smoothed state estimates. In object tracking, a common state formulation is \([x, y, v_x, v_y]\) (position and velocity). The filter alternates between prediction (using a motion model) and update (incorporating measurements), which makes it effective for smoothing noisy trajectories and interpolating short gaps in detections.

In the Padel Trainer pipeline, a constant-velocity Kalman Filter is used to smooth the ball trajectory and to provide short-term predictions when detections are missing.

\subsection{Edge AI}

Edge AI performs inference on local or embedded devices rather than in the cloud, reducing latency, bandwidth, and privacy risks. Edge deployment is beneficial for real-time sports analytics because it enables immediate feedback and offline operation. However, resource constraints on edge hardware (CPU, memory, power) require model and system optimization.

\subsubsection{Raspberry Pi and Hailo Accelerator}

The Raspberry Pi 5 is an ARM-based single-board computer suitable for embedded workloads. For demanding deep learning inference, external accelerators such as the Hailo-8 can provide substantial throughput improvements while maintaining low power consumption. The target edge platform for this work combines a Raspberry Pi 5 with a Hailo accelerator to meet real-time processing requirements on the court.

\subsection{ONNX}

The Open Neural Network Exchange (ONNX) is an open format for representing machine learning models that enables interoperability across frameworks and optimized runtimes. Exporting models (for example from PyTorch) to ONNX simplifies deployment on heterogeneous hardware and allows using optimized inference engines on edge devices.